
\chapter{Санкционный скрининг}\label{ch:sanctions}
Санкционный скрининг -- отдельный контур контроля, не сводимый
к~AML. У~него свои перечни, свои алгоритмы сравнения, свои пороги
ложных срабатываний и свой маршрут эскалации. Российский оператор
работает с~тремя национальными перечнями; международные перечни
(OFAC, EU, UN) применяются по собственной риск-политике
и~требованиям банков-корреспондентов. Если клиент, бенефициар
или контрагент попадает в~применимый перечень, операцию
останавливают до исполнения.

\section{Российские национальные перечни}\label{sec:sanctions-rf-lists}

Для российского оператора первичны три национальных перечня.
Во-первых, перечень организаций и физических лиц, в~отношении
которых имеются сведения об их причастности к~экстремистской
деятельности или терроризму, который ведёт Росфинмониторинг;
меры по замораживанию средств таких лиц применяются по
подпункту~6 пункта~1 статьи~7 115-ФЗ (операции с~лицами
из перечня дополнительно подлежат обязательному контролю
по статье~6). Во-вторых, перечень организаций и физических лиц,
в~отношении которых имеются сведения об их причастности
к~распространению оружия массового уничтожения. В-третьих,
решения Межведомственной комиссии (МВК) о~замораживании
(блокировании) денежных средств или иного
имущества~\cite{rosfinmonitoring-lists,fz-115,cbr-mvk}.\footnote{Точные
наименования перечней, формат публикации и периодичность
обновления сверяются по сайту Росфинмониторинга и по
действующей редакции 115-ФЗ.}

Перечни Росфинмониторинга обновляются регулярно; после публикации
обновлённого перечня оператор обязан в~установленный срок не позднее
одного рабочего дня прогнать обновление и применить меры к~клиентам,
попавшим в~перечень~\cite{rosfinmonitoring-lists,fz-115}.\footnote{Срок
применения мер к~клиентам после обновления перечня -- предмет
последовательных правок 115-ФЗ; перед публикацией сверить
по действующей редакции.}

\section{Международные перечни и их статус для российского
оператора}\label{sec:sanctions-international}

Перечни ООН, ЕС и OFAC SDN для российского оператора имеют статус,
отличный от национальных. Их использование определяется собственной
риск-политикой оператора, требованиями банков-корреспондентов
и~санкционным режимом конкретной транзакции; они не подменяют
национальные
перечни~\cite{ofac-sanctions-list-service,eu-sanctions-info,unsc-consolidated-list}.

На практике это означает следующее. Транзакция в~валюте
недружественной страны почти неизбежно проходит через
банк-корреспондент, который применяет соответствующие санкционные
режимы; даже если российский оператор формально не обязан
проверять контрагента по OFAC SDN, отказ корреспондента
в~исполнении приведёт к~операционной блокировке. Соответственно,
оператор включает международные перечни в~собственный скрининг
не потому, что этого требует национальный закон, а~ради
устойчивости расчётного канала.

\section{Архитектура скрининга}\label{sec:sanctions-screening-flow}

На практике оператор ставит проверку по перечням как минимум
на подключении, при операциях, попадающих в~риск-профиль,
при обновлении данных клиента и~при обновлении самих перечней.

Скрининг одной транзакции проходит так. Авторизационный запрос
поступает на хост эмитента. Ещё до проверки лимитов модуль
скрининга извлекает имя держателя из CMS по PAN, нормализует его
(транслитерация, удаление диакритики, разбиение на n-граммы)
и сравнивает с~индексом санкционных перечней в~оперативной
памяти. Если нечёткое совпадение превышает порог, транзакция
уходит в~очередь ручной проверки; оператор видит контекст
(страна, сумма, история клиента) и принимает решение: снять блок
или эскалировать. Если совпадение ниже порога, транзакция
продолжает обычный путь авторизации; добавленная латентность
ограничена временем lookup в~индексе.\footnote{Конкретные
числовые пороги (доля совпадения по Jaro--Winkler и~пр.)
и~время разбора кейса в~книге не приводятся: они зависят
от настройки конкретного оператора и не зафиксированы в
нормативных актах.}

Скрининг в~онлайне тормозит авторизацию, и~бизнес ищет компромисс
между скоростью и полнотой. Рабочая система строится на индексах
в~оперативной памяти и~на чёткой очереди эскалации: ручная проверка
каждого совпадения не масштабируется. Распространённые имена
и совпадения между кириллицей и латиницей порождают массу шума;
задача -- находить подозрительные совпадения и быстро разбирать
ложные.

\section{Ложные срабатывания и маршрут разбора}\label{sec:sanctions-false-positives}

Ложноположительные совпадения требуют отдельного маршрута разбора.
Типовой сценарий: система помечает операцию как требующую проверки,
сотрудник смотрит контекст, после чего либо снимает блок, либо
эскалирует случай. Маршрут должен быть формализован, иначе итог
зависит от того, кто дежурит на смене. Слишком мягкий порог топит
команду в~шуме; слишком жёсткий пропускает риск. Логика сопоставления
и пороги срабатывания настраиваются исходя из собственной оценки
риска оператора, зафиксированной в~правилах внутреннего контроля
по 115-ФЗ.
